%! TEX root = formfollowfunction.tex
% the main TeX file which is intended to compile, :VimtexReload after adjustment
   % this file should be included in the main file
% :h vimtex-tex-root
\documentclass[../manperil.tex]{subfiles}
%ensure the class of docs

\begin{document}


\section{Fungsi}

Christian Norberg Schulz menulis sebuah buku yang berjudul Intention in Architecture yangintention itu sendiri dapat diartikan sebagai makna, tujuan, kegunaan, atau umumnya disebut sebagai fungsi dalam arsitektur. Dalam bukunya, ia mengatakan bahwa semua hal pengetahuan teori maupun praktek yangberhubungan dengan fungsi arsitektur berawal dari pertanyaan “What is architecture?” atau “Apakaharsitektur itu?”. Menurut C. Norberg Schulz, arsitektur merupakan suatu lingkungan yang dibuat atau dibina oleh manusia (Built environment).

Namun tidak hanya sampai pada pertanyaan itu, kemudian muncul lagi pertanyaan lanjutan yangdijadikan landasan tahap pengembangan teori “fungsi” arsitektur ini. Norberg-Schulz menyatakan ketigakategori dari fungsi arsitektur, yaitu “functional-practical purposes, milieu-creating purposes, symbolizingpurposes”. Seperti Broadbent, fungsi arsitektur ini disebut “Building Task”. Schulz yakin bahwa fungsiarsitektur tidak dapat dilihat dari permukaan fisiknya saja, tapi juga secara sosiokultural. Maka dari itu, iamenggabungkan teori fungsi arsitekturnya dengan mengajukan beberapa kategori fungsi arsitektur yangdisampaikan dalam 4 poin, yaitu:

\subsection{Macam-macam Fungsi}%
\label{sub:Macam-macam Fungsi}


\subsubsection{Physical Control (Pengendali Alam)}%
\label{sub:physical_control_(pengendali_alam)}


Mengenai fungsi dan peran sebuah bangunan dalam mengendalikan kondisi alam sepertimelindungi manusia dari masalah pergantian cuaca, melindungi dari bencana, dan lainnya. Karena itu,arsitek membutuhkan abstraksi tentang segala sesuatu yang berhubungan langsung dengan aspek fisikbangunan tersebut. Misalnya, dapat diselidiki dari kemampuan bahan bagunan sebagai insulator terhadapdingin, suara, kelembaban, dll. Dapat juga menggunakan alat-alat yang secara mekanis untuk menghasilkan“iklim artifisial”. Fungsi bangunan sebagai pengendali ala mini berkaitan erat dengan konsep “greenbuilding”. Contohnya pada bangunan NTU (Nanyang School of Art) yang merupakan bangunan dengankonsep green architecture terbaik di Singapore dan juga Reichstag New German Parliament yang di desainpada puncak gedung parlemen dengan konsep merespon alam.


\subsubsection{Functional Frame (Kerangka Fungsi)}%
\label{sub:functional_frame_(kerangka_fungsi)}

Arsitektur dapat menciptakan kerangka fungsi, atau bisa dibilang sebagai bangunan atau ruanganyang dibedakan dari aktivitas yang terjadi didalamnya. Umumnya manusia memang selalu beraktivitas,sehingga membutuhkan wadah untuk menampung aktivitas tersebut. Dua bangunan dapat berperan denganbaik untuk fungsi yang sama meski suasananya berbeda. Kalau mengenai aktivitas yang menjadikan fungsi,maka hal ini sangat berkaitan dengan kajian tipologi bangunan. Contohnya seperti gereja dan masjid, keduabangunan yang fungsinya sama-sama untuk beribadah namun suasananya berbeda.
Gereja Katedral Jakarta memiliki nuansa yang diadaptasi dari Eropa, sedangkan Masjid Istiqlalmengadaptasi budaya Timur Tengah untuk bentuk bangunannya.


\subsubsection{Social Milieu (Lingkungan Sosial)}%
\label{sub:social_milieu_(lingkungan_sosial)}


Bangunan mampu membentuk lingkungan sosial sendiri. Arsitektur sering dianggap sebagai alatuntuk mengekspresikan status sosial pemilik atau penggunanya. Ini merupakan bukti bahwa arsitekturdiekspresikan untuk maksud dan tujuan tertentu khususnya dalam tingkat sosial masyarakat setempat. Jugaarsitektur diharapkan untuk menjamin proses interaksi sosial berlangsung dengan sebagaimana mestinyadan memberi pengaruh psikologis terhadap lingkungan sosial budayanya. Contohnya rumah adat Tongkonan yang ada di Tana Toraja. Pada ornamen kerbau yang ada di depan rumah, semakin banyakkepala kerbau yang dipasang, menandakan semakin tinggi derajat sosial pemilk rumah tersebut dilingkungan tersebut. Juga rumah yang ada di Eropa, biasanya semakin luas taman halaman rumahnya,berarti semakin kaya juga si pemilik.

\subsubsection{Cultural Symbolization (Simbol Budaya)}%
\label{sub:cultural_symbolization_(simbol_budaya)}

Bangunan dapat menjadi simbol budaya masyarakat setempat. Arsitektur merupakan sebuah objekbudaya yang juga merupakan hasil karya manusia untuk melayani aktivitas manusia secara umum. Denganberbagai olahan, arsitektur senantiasa mengekspresikan nilai dan sistem sosial budaya masyarakat dimanaarsitektur itu di dirikan. Contohnya Hagia Sophia yang ada di Istanbul, Turki yang merupakan salah satubangunan terkenal mengisyaratkan bahwa Turki kental dengan budaya timur tengah dengan konsepbangunan dengan dome dan menara untuk rumah ibadah. Dan juga seperti rumah-rumah yang ada di Bali,Indonesia yang masih menerapkan budaya mereka dalam membangun sebuah rumah tinggal yang bisadisebut juga compound yang membuat wisatawan sadar akan budaya yang dimiliki Bali.

\subsection{Prinsip Keseimbangan Desain dan Fungsionalitas}%
\label{sec:Prinsip Keseimbangan Desain dan Fungsionalitas}

\subsubsection{Desain Berpusat pada Pengguna}%
\label{sub:desain_berpusat_pada_pengguna}

Prinsip ini merupakan inti dari seluruh pedoman yang ada dan harus selalu dipertahankan sebagai landasan utama. Konsep dasar ini menitikberatkan pada keharusan untuk memahami kebutuhan karyawan dan melibatkan mereka dengan penuh empati. Proses ini mencakup identifikasi komprehensif terhadap berbagai aspek pengguna ruang, antara lain:

\begin{itemize}
	\item Kebutuhan spesifik
	\item Kebiasaan kerja
	\item Preferensi desain
	\item Perilaku pengguna
\end{itemize}


\begin{itemize}
	\item Pelaksanaan Survei: Mengumpulkan data terukur mengenai kebutuhan dan harapan pengguna.
	\item Penyelenggaraan Lokakarya dan Dialog: Memfasilitasi diskusi dua arah untuk menggali wawasan secara lebih mendalam.
	\item Pembuatan Ruang Percontohan (Pilot Space): Menyediakan area uji coba untuk mengobservasi secara langsung bagaimana pengguna berinteraksi dengan lingkungan kerja mereka.
\end{itemize}


\subsubsection{Fungsionalitas sebagai fondasi}%
\label{sub:fungsionalitas_sebagai_fondasi}

Fungsionalitas merupakan landasan utama dalam setiap pengembangan tata ruang. Meskipun estetika visual memiliki nilai penting, prioritas utama dari sebuah ruang kerja adalah kemampuannya dalam mendukung pelaksanaan tugas dan aktivitas operasional secara optimal. Penerapan standar fungsionalitas ini mencakup:

\begin{enumerate}
	\item Penyediaan area kerja (workstation) yang memenuhi standar ergonomis.
	\item Integrasi infrastruktur teknologi yang mulus dan menyeluruh.
	\item Kemudahan mobilitas dan navigasi di dalam ruang kerja.
\end{enumerate}


Desain kantor yang optimal harus memiliki sifat intuitif bagi para penggunanya. Untuk menciptakan lingkungan kerja yang efisien dan ramah pengguna (user-friendly), terdapat beberapa elemen tata ruang kunci yang harus diimplementasikan:

\begin{enumerate}
	\item Sistem Penunjuk Arah (Wayfinding): Pemasangan rambu atau papan informasi penunjuk arah yang jelas dan informatif.
	\item Tata Letak (Layout) Logis: Pengaturan alur ruang dan penempatan zona kerja yang terstruktur secara sistematis.
	\item Aksesibilitas Fasilitas: Penyediaan berbagai fasilitas penunjang yang mudah dijangkau dan dapat diakses oleh seluruh karyawan.
\end{enumerate}

\subsubsection{Estetika sebagai elemen harmonisasi}%
\label{sub:estetika_sebagai_elemen_harmonisasi}


Meskipun fungsionalitas merupakan landasan dasar, desain dan estetika memegang peranan krusial dalam menciptakan lingkungan kerja yang menginspirasi dan menyenangkan. Elemen estetika harus diintegrasikan secara strategis untuk melengkapi aspek fungsional, sehingga tercipta tata ruang yang kohesif dan harmonis.

Optimalisasi suasana ruang secara keseluruhan dapat dicapai melalui penerapan elemen-elemen visual berikut:

\begin{enumerate}
	\item Skema Warna: Pemilihan palet warna yang mendukung dinamika kerja dan psikologi ruang.
	\item Material dan Tekstur: Penggunaan bahan yang memberikan nilai visual sekaligus kenyamanan taktil.
	\item Sistem Pencahayaan: Tata cahaya yang dirancang untuk memaksimalkan fungsi sekaligus meningkatkan kualitas atmosfer ruangan.
\end{enumerate}

\subsubsection{Keberlanjutan}%
\label{sub:Keberlanjutan}

Dalam konteks pengembangan ruang kerja modern, penerapan desain berkelanjutan bukan lagi sekadar pilihan, melainkan sebuah standar keharusan yang menjadi prinsip inti tata ruang. Implementasi konsep ramah lingkungan ini mencakup beberapa indikator utama:

\begin{enumerate}
	\item Penggunaan material bangunan dan furnitur yang ramah lingkungan (eco-friendly).
	\item Pemasangan sistem utilitas dan operasional yang hemat energi.
	\item Perancangan struktur yang difokuskan pada pengurangan jejak karbon (carbon footprint) bangunan.
\end{enumerate}

\subsubsection{Dampak Optimalisasi Keberlanjutan}%
\label{ssub:dampak_optimalisasi_keberlanjutan}

Integrasi prinsip keberlanjutan secara langsung akan meningkatkan kualitas desain sekaligus fungsionalitas ruang kerja. Penerapan pendekatan ini dirancang untuk memberikan keuntungan ganda, yang meliputi:

\begin{enumerate}
	\item Peningkatan Kesehatan Lingkungan: Menciptakan ruang kerja yang lebih sehat untuk mendukung kesejahteraan dan produktivitas pengguna.
	\item Efisiensi Finansial Jangka Panjang: Menekan beban biaya operasional dan pemeliharaan gedung secara signifikan di masa mendatang.
\end{enumerate}

\subsubsection{Viabilitas Jangka Panjang dan Keselarasan ESG}%
\label{ssub:viabilitas_jangka_panjang_dan_keselarasan_esg}

Pendekatan tata ruang berbasis keberlanjutan berfungsi sebagai jaminan bahwa ruang kerja akan tetap layak, adaptif, dan efisien dalam jangka waktu yang panjang. Lebih jauh, strategi spasial ini memastikan bahwa aset fisik perusahaan sejalan dengan visi makro organisasi, yaitu:

\begin{enumerate}
	\item Pemenuhan komitmen Tanggung Jawab Sosial Perusahaan (Corporate Social Responsibility/CSR).
	\item Pencapaian target dan metrik Environmental, Social, and Governance (ESG) yang telah ditetapkan.
\end{enumerate}




Sumber : \textit{https://id.linkedin.com/pulse/core-principles-guiding-balance-between-design-bharat-tagra--ke7oc?tl=id}

\subsection{Macam-macam Fungsi dalam Arsitektur}%
\label{sub:macam-macam_fungsi_dalam_arsitektur}

Fungsi adalah tugas dan pekerjaan yang harusdijalankan oleh sebuah lingkungan.maka fungsi dalam arsitektur dapat dikelompokkan menjadi:

\subsubsection{Fungsionalisme Bentuk}%
\label{ssub:fungsionalisme_bentuk}


Fungsionalisme bentuk merupakan fungsi yang paling lazim dan paling mudah ditangkapdari sebuah produk arsitektur atau lingkungan hasil binaan manusia. Karena fungsi inimemiliki makna paling awal; Form Follow Function (bentuk mengikuti fungsi), dimanasegala rancangan arsitektur terjadi karena fungsi, pembedaan bagian bangunan menuruttujuannya, rancangan bangunan untuk memenuhi kebutuhan manusia dan bentuk berasal darikeinginan manuasia yang akan menggunakannya sesuai aktivitasnya.

\subsubsection{Fungsionalisme Konstruksi}%
\label{ssub:fungsionalisme_konstruksi}


Struktur, konstruksi dan bahan bangunan sampai batas tertentu memiliki kedudukan yanglebih tinggi, Form Follow Structure Function (bentuk terjadi akibat syarat-syarat konstruksi),dimana bentuk berasal dari syarat sistem struktur, konstruksi dan bahan bangunan menurutpenggunaan struktur, konstruksi yang jujur, jelas dan wajar tanpa disembunyikan, rancanganstruktur untuk tujuan estetik yang terwujud melalui elemen-elemen struktur itu sendiri.Contoh dari fungsi ini adalah menara Burjj Khalifa, Dubai. Menara tersebut memilikipondasi segitiga, dimana pondasi ini merupakan pondasi yang paling ideal digunakan untukbangunan-bangunan pencakar langit. Menara yang mempunyai berat 500.000 ton ini berdiridi atas pondasi berbentuk segitiga dengan ketebalan 3,7m. Didukung oleh rangka-rangkabaja yang mengelilingi dan mengisi pondasi dengan diameter masing 1,5m. Pondasi inimenyebabkan terciptanya bentuk burj khalifa yang terdiri dari tiga elemen atau sayap yangdisusun mengelilingi satu pusat yang berada tepat di tengah. Struktur yang berbentuk Ymemberikan kestabilan pada masing-masing (ketiga) elemen/sayap, dimana struktur tersebutmerupakan landasan yang baik dan ideal untuk pembuatan residen. Dengan struktur Y,memaksimalkan pemandangan ke arah Teluk Arab. Burj Khalifa bila dilihat dari atasmenyerupai kubah masjid seperti lazimnya bangunan-bangunan berarsitektur Islam.Dengan demikian, struktur dari Burj Khalifa mempengaruhi bentuk dari Burj Khalifa.

\subsubsection{Fungsionalisme Ekspresi}%
\label{ssub:fungsionalisme_ekspresi}


Fungsi ekspresi memperlihatkan guna dan struktur secara bersama-sama dalam arsitektur,dimana bentuk merupakan wujud dari kegunaan / fungsi di dalamnya, bentuk secarasimbolik melukiskan fungsi, rancangan bangunan memperlihatkan struktur \& konstruksiserta peralatan bangunan secara menonjolSebuah rumah di Korea Selatan memiliki bentuk yang unik dan menarik. Setiap orang pastiakan dengan sangat mudah mengenali bantuk yang diadopsi oleh bangunan seluas 419 m2ini. Bangunan bernama Haewooajae ( yang dalam bahasa Korea “berarti sebuah tempat yangmemberikan solusi kecemasan bagi seseorang” ) sengaja dibuat identik dengan bentuk klosetduduk, agar sesuai dengan kegunaan dominan yang dimiliki bangunan tersebut.Kaitannya dengan Toilet House, produk arsitektur ini menyajikan makna fungsi ekspresi. Toilet House yang berbentuk kloset yang merupakan bagian dari toilet adalah bangunan yangmempunyai fungsi utama untuk sanitasi. Sang pemilik, Sim Jae-duck yang merupakan salahsatu orang berpengaruh di Asosiasi Toilet Dunia (World Toilet Association) inginmempertegas fungsi dari bangunan ini dengan mengambil bentuk dari elemen toilet yangsimbolik, yaitu kloset duduk. Dengan demikian, bentuk dan struktur yang ditampilkanselaras dengan kegunaan bangunan tersebut, yaitu berbentuk kloset, dan memiliki kegunaanutama untuk sanitasi.Bagian luar bangunan ini dilapisi baja yang dicat putih mirip dengan permukaan klosetkeramik. Terdapat juga kaca-kaca tidak tembus pandang sebagai pelapis dinding luar.Bukaan di atas bangunan yang serupa dengan lubang kloset, murupakan pintu masuk dariToilet House ini.Toilet House terdiri dari dua lantai, berisi tiga kamar mandi mewah dan komplit, serta 2 kamar tidur. Dinding dalam bangunan dan jendela-jendela lengkung cantik juga mengikutipermukaan bangunan. Nuansa yang ditampilkan adalah nuansa bersih, penggunaan catkeramik putih sebagai pelapis dinding bagian dalam memperkuat citra kloset pada bangunanini.Banguan ini dilengkapi dengan tombol-tombol sensor otomatis. Seperti ketika seseorangmemasuki kamar mandi ini, tombol sensor otomatis akan memainkan musik-musik klasikyang diharapkan membuat nyaman pengguna. Kini, Sim Jae-duck yang selalumengkampanyekan untuk mempercantik toilet yang higienis ini tengah mencari seseorangyang mau menginap semalam di rumah toilet dengan membayar usd50.000. Dana inirencananya akan digunakannya untuk membantu program pemerintah dalam rangkamenciptakan toilet yang sehat dan bersih di negaranya.Toilet House ini membuktikan bahwa produk arsitektur memiliki makna fungsi espresif,dimana kegunaan bangunan dapat tercermin dan tersimbolkan oleh struktur dan bentukan-bentukan ekspresif suatu bangunan.


\subsubsection{Fungsionalisme Geometris}%
\label{sub:fungsionalisme_geometris}


Fungsi ini mencoba mengabaikan kegunaan bangunan dan memusatkan perhatian pada caradimana geometri bangunan berfungsi secara visual dan estetis.

Kebalikan dari fungsi bentuk, prinsip Function Follow Form mengiringi fungsi geometri ini,dimana penciptaan bentuk bukan untuk menyesuaikan dengan guna, tetapi akibatpenyesuaian bentuk geometris itu sendiri, kesederhanaan bentuk dengan geometri dan bebasdari ornament, nilai estetis didapat dari pengolahan elemen geometri.Contoh dari fungsi ini adalah Wooden House karya Sou Fujimoto di Kumamoto, Jepang.Dalam wooden house yang didesain dari tahun 2006 sampai tahun 2006 dan berbentuk kubusini, balok-balok kayu ditata sedemikian rupa sehingga menciptakan produk arsitektur yanggeometris. Sedangkan aktivitas yang terjadi di dalamnya mengikuti bentuk dari woodenhouse tersebut. Balok kayu yang terbentang dengan permukaan lebar dan memanjang sesuaiuntuk aktivitas berbaring atau tidur, sedangkan balok yang mempunyai permukaanhorizontal sempit namun mempunyai permukaan vertical tinggi bisa digunakan untuk duduklayaknya kursi. Di sinii terlihat bahwa fungsi menyesuaikan atau dipengaruhi oleh bentukbangunan itu sendiri.

\subsubsection{Fungsionalisme Organis}%
\label{ssub:fungsionalisme_organis}


Karya arsitektur tidak hanya fungsional tetapi juga organis (bentuk terinspirasi darikehidupan yang alamiah), dimana bentuk dan fungsi identik, karya arsitektur berwawasanlingkungan, bentuk terinspirasi fenomena alam dan penggalian gagasan dari mahluk hidup,fungsi bangunan adalah aktifitas yang menciptaka bentuk, sehingga bentuk adalah fungsidari keseluruhan.Arsitektur yang dipahami sebagai seni dan keseluruhan proses membangun, ternyata bukanhanya dilakukan oleh manusia, tetapi juga oleh alam seisinya. Binatang-binatang yangdengan instingnya –sendiri maupun bersama koloninya- membangun sebuah produkarsitektur untuk memenuhi kebutuhan hidupnya, entah itu untuk tempat tinggal, tempatmencari mangsa, atau pun untuk tujuan lainnya sebagai wujud dari eksistensi dirinya, disebutdengan arsitektur alam (natural architecture). Bahkan binatang-binatang sudah terlebihdahulu bisa ‘hidup’ di alam ini, secara alamiah telah membangun sarang-salah satunya-dengan kompleksitas melebihi produk arsitektur manusia yang baru ditemukan berabad-abadkemudian.Manusia sebagai makhluk yang mempunyai kebutuhan akan ruang sebagai tempatmelangsungkan kehidupannya dengan mapan dan kerasan perlulah manusia -sendirian ataubersama dengan sesamanya- berada dalam kesesuaian dengan alam dan seisinya. Bisa jadi,dialah yang harus menyesuaikan diri dengan keadaan alam sekitarnya (dinamakan tindakanakomodasi), namun bisa jadi pula sebaliknya: ia mengubah alam agar sesuai dengan polakehidupannya (tindakan asimilasi). Sepanjang hidupnya manusia harus mengusahakankeseimbangan antara akomodasi dan asimilasi itu, yang disebut sebagai tindakan adaptasi.Adaptasi itulah yang hendak kita usahakan dalam ber-arsitektur: yaitu tercapainyakeselarasan manusia dengan alam melalui gubahan massa dan ruang. Membangun arsitektur-dengan demikian- adalah berarti menciptakan suatu lingkungan yang di dalamnya berisi(menampung) keselarasan antara pola kehidupan manusia penghuninya dengan alam diluarnya.Yang pertama kali terlintas dalam benak seseorang ketika berpikir tentang laba-laba adalahjaringnya. Ia merupakan keajaiban desain yang memiliki rancangan tersendiri, besertaperhitungan teknik yang menyertainya. Jika kita memperbesar laba-laba menjadi seukuranmanusia, jaring yang dianyamnya akan memiliki tinggi sekitar seratus lima puluh meter. Inisama tingginya dengan gedung pencakar langit berlantai lima puluh.Andaikan laba-laba sedemikian besar sehingga mampu membuat jaring dengan lebar limapuluh meter, maka jaring ini akan mampu menghentikan pesawat jumbo jet. Jika demikian,bagaimana laba-laba mampu membuat jaring dengan sifat ini? Agar dapat melakukan hal ini,ia pertama kali harus menggambar rancangannya, persis seperti seorang arsitek. Sebab,struktur arsitektural dengan ukuran dan kekuatan seperti ini, mustahil dilakukan tanpa sebuahperancangan. Setelah rancangan dipersiapkan, laba-laba perlu menghitung seberapa besarbeban-beban yang akan menempati posisi-posisi tertentu pada jaring, persis layaknyainsinyur konstruksi. Jika tidak, jaring ini pasti akan runtuh. Namun kenyataannya, laba-labamembangun itu semua tanpa perancangan, hanya insting dari Sang Pencipta lah yangmenuntunnya hingga mampu membuat sarang sedemikian lengkap dengan unsur kekuatan,elastisitas maupun keindahan.Contoh lain dari fungsi organis adalah bisa kita lihat sekali lagi pada Burj Khalifa. Biladipandang dari udara kita juga bisa melihat bahwa tampak atas Burj Khalifa menyerupaikelopak bunga; Hymenocallis, dari batang hingga sayap menara yang membentang dari intipusatnya. Arsitek Adrian Smith sedikit banyak memasukkan pola dari arsitektur Islamtradisional, tetapi bentuk dan konsep yang paling inspiratif dan menonjol adalah bunga gurunlokal, yaitu Hymenocallis.Desain bunga Hymenocallis tidak hanya mampu mengurangi tekanan angin yang akanmengarah kepada bangunan tapi juga menjadikan setiap sisi mempunyai pemandangan TelukArab yang luar biasa.


\subsubsection{Fungsionalisme Ekonomis}%
\label{ssub:fungsionalisme_ekonomis}


Menggunakan pendekatan ekonomi dalam proses penciptaan karya Arsitektur, dimanabentuk terjadi akibat pemakaian peralatan dan bahan secara ekonomis, penggunaan metodedan cara yang paling efektif dan efisien agar bisa menyesuaikan biaya yang dikeluarkandalam proses pembuatan dengan biaya yang disediakan.Contoh dari fungsi ini adalah kembali lagi kita melihat Wooden House karya Sou Fujimotodi Kumamoto, Jepang. Sou Fujimoto menggunakan satu bahan bangunan yang dominanuntuk membangun seluruh bagian dari Wooden House ini. Fungsi ekonomis terjadi di sini,dimana harga kayu akan lebih ekonomis bila dibandingkan dengan bahan bangunan laindengan mangingat bahwa kayu memiliki ketahanan yang sangat kuat terhadap fenomenaalam seperti gempa, sehingga fengan masa ketahanan yang lama, akan mamperkecil biaya pembuatan serta perawatan.

\subsubsection{Fungsionalisme Kultural}%
\label{ssub:fungsionalisme_kultural}

Fungsi ini menempatkan manusia secara sentral sebagai pusat, pedoman dalam menciptakanbentuk karya arsitektur. Mengikuti prinsip Form Follow Culture, dimana bentuk berasal daripola perilaku, tradisi, adat istiadat, kondisi sosial budaya pemakai, bentuk dijiwai olehkehidupan manusia, watak, kecenderungan dan nafsu serta cita-cita.Manusia dalam berarsitektur hendaknya memahami bahwa faktor-faktor yang mempengaruhibahkan penentu produk arsitekturnya berkaitan erat dengan alam dan seisinya, diantaranya;tradisi masyarakatnya sendiri, pola kegiatan masyarakat dan tentunya faktor alam, seperti keadaan geografis suatu wilayah -perlu kita ketahui ciri-sifatnya (karakternya)- untukmengetahui bagaimana teknologi (bahan,alat dan metode) yang dipilihnya melahirkan bentukarsitektur tertentu, dan iklim dari lokasi produk arsitektur.Rumah masyarakat Jawa Tengah misalnya, yang mempunyai tradisi mengagungkan NyiRoro Kidul, mereka membangun rumah menghadap selatan sebagai penghormatan terhadapRatu Pantai Selatan, ini menunjukkan bahwa factor tradisi masyarakat mempengaruhiarsitektur mereka.Faktor geografis yang berpengaruh terhadap arsitektur bisa dilihat dari rumah panggung diMinahasa misalnya, mereka menyesuaikan diri dengan alam sekitar dimana terdapat banyakbinatang buas dengan membuat rumah yang berbentuk panggung.

\subsection{Penerapan Fungsi pada Bangunan}%
\label{sub:penerapan_fungsi_pada_bangunan}


\subsubsection{Artistik Form (Fungsi Bentuk Arsitektur)}%
\label{ssub:artistik_form_(fungsi_bentuk_arsitektur)}

Pada entrance diekspose kolom-kolom bulat selainberfungsi sebagai struktur berfungsi juga sebagaiestetika. Pada bagian badan bangunan lebih menekankanpada keseimbangan bangunan, baik itu keseimbanganbentuk, grid dan struktur yang kesemuanya nampak padatampilan bangunan dan unsur kajian tipologi objekterhadap bangunan lain.

Bentuk arstistic pada bangunan yang memiliki keindahan tersendiri untuk memikat para pemakai, penikmat dan pengamat.

\subsubsection{Container (FungsiPerwadahan)}%
\label{ssub:Container (FungsiPerwadahan)}

Dibuat hal atau lobby sebagai perwadahan aktifitas.

\subsubsection{Climatic Modifier(Fungsi Modifikasi /Kontrol Iklim)}%
\label{ssub:climatic_modifier(fungsi_modifikasi_/kontrol_iklim)}

Penggunaan material, khusunya material selubungbangunan khususnya pada skylight hall penerima utamadan hall terbuka untuk mendapat pencahayaan alami.· Penggunaan barrier dalam menanggapi masalahyang timbul karena angin dan hujan· Penggunaan drainase yang baik dan terjaga, jugapenggunaan elemen-elemen yang dapat mengantisipasibanjir akibat kapasitas air hujan yang besar.

\subsubsection{Environmental Filter (Fungsi Filter Lingkungan}%
\label{sub:Environmental Filter (Fungsi Filter Lingkungan}

Suasana dibuat dapat memberikan suatu daya tarik bagi lingkungan sekitar sehingga memberi kesanmenyatu dengan bangunan sekitar.

\subsubsection{Behaviour Modifier(Fungsi Pembentuk Perilaku)}%
\label{sub:behaviour_modifier(fungsi_pembentuk_perilaku)}

Penggunaan material pada dinding, kolom, danbalok berbeda pada setiap ruangan menjadikan fungsipembentuk perilaku.· Disekitar pedestrian way ditanami tumbuhansepanjang jalan dan lampu jalan sebagaimana fungsipengantar pengunjung menuju pintu masuk.· Penataan, hubungan dan dimensi ruang pada setiapfasilitasnya sehingga akan melekat pada pengunjungyang datang diHotel Bisnis.

\subsubsection{Capital Investment(Fungsi InvestasiModal)}%
\label{sub:capital_investment(fungsi_investasimodal)}
Hotel Bisnis ini bertujuan untuk mencari keuntungan.Keuntungan yang diinvestasikan melalui objek iniberupa keuntungan profit dan benefit. Keuntungan profitdidapat melalui pemakaian fungsi ruang yang ada sertajumlah pengunjung. Sedangkan keuntungan benefitdiperoleh melalui manfaat yang didapat daripengunjung/pemakai akan fasilitas yang ada pada Hotel Bisnis ini.

\subsubsection{Cultural Symbolization (Fungsi SimbolBudaya)}%
\label{sub:cultural_symbolization_(fungsi_simbolbudaya)}

Memberikan simbol budaya pada bangunan dengan memasukkan unsur-unsur yang berkaitan dengan budayadaerah dimana objek berada.

\section{Estetika}%
\label{sec:Estetika}

\subsection{Pengertian Estetika}%
\label{sub:Pengertian Estetika}
Secara etimologis, istilah “estetika” berasal dari bahasaLatin “aestheticus” atau bahasa Yunani “aestheticos” yang artinya merasa atauhal-hal yang dapat dicerap oleh panca indera manusia. Ada juga yangmenyebutkan bahwa arti estetika adalah suatu cabang ilmu filsafat yangmembahas tentang keindahan, dan biasanya terdapat di dalam seni dan alamsemesta. Sedangkan menurut KBBI, istilah estetika memiliki dua arti. Estetikaadalah suatu cabang filsafat yang membahas tentang seni, nilai keindahan, dantanggapan manusia terhadapnya. Estetika juga dapat didefinisikan sebagaikepekaan manusia terhadap seni dan keindahan.

Dari penjelasan tersebut dapat disimpulkan bahwa estetika sangatberkaitan dengan perasaan manusia, khususnya perasaan yang indah atau perasaanpositif. Keindahan yang dimaksud di sini bukan hanya sesuatu yang dapat dilihatbentuknya, tapi juga makna atau arti yang terkandung di dalamnya.Sedangkan estetika menurut beberapa ahli ialah:

\begin{enumerate}
	\item Bruce Allsopp
	\item[] Menurut Bruce Allsopp (1977), estetika adalah ilmu pengetahuanyang mempelajari tentang proses dan aturan dalam menciptakan suatu karya seni, yang diharapkan bisa menimbulkan perasaan positif bagi orangyang melihat dan merasakannya.

	\item Herbert Read
	\item[] Menurut Herbert Read, pengertian estetika adalah kesatuan danhubungan bentuk yang ada di antara pencerapan indrawi manusia.Biasanya, manusia menganggap estetika adalah seni, atau seni akan selalumengandung nilai keindahan. Pandangan seperti ini sebenarnya salah danjustru membuat masyarakat kesulitan dalam mengapresiasi seni. Sebab,menurut Herbert Read, seni tidak selalu mengandung nilai keindahan.

	\item J. W. MorisA
	\item[] Menurut J. W. Moris, pengertian estetika sama dengan seni karenaestetika dapat dikenakan pada berbagai objek, baik yang indah maupuntidak. Selanjutnya, Moris juga menyebutkan bahwa estetikan adalah suatuobjek seni (art).

\end{enumerate}

Terdapat beberapa unsur yang perlu diperhatikan dalam mengkaji nilaiestetika suatu objek. Adapaun beberapa unsur estetika adalah sebagai berikut:


\begin{enumerate}

	\item Unsur Bentuk
	\item[] Bentuk (shape) sangat berpengaruh pada daya tarik suatu objek.Secara umum, bentuk objek terdiri dari dua jenis, yaitu; dua dimensi dantiga dimensi. Objek berbentuk dua dimensi tidak memiliki volume danbentuknya datar. Misalnya lukisan, foto, hiasan dinding, dan lainnya. Objek berbentuk tiga dimensi memiliki volume, kedalaman, dan ruang.Misalnya patung, pakaian, tas, dan lainnya.

	\item Unsur Warna
	\item[] Keindahan suatu objek juga sangat dipengaruhi oleh unsur warna.Umumnya pilihan warna objek akan disesuaikan oleh orang yang akanmenggunakannya. Misalnya, selera warna pakaian anak muda cenderungberbeda dengan orang yang sudah tua.


	\item Unsur Tema
	\item[] Dalam hal ini tema adalah ide atau gagasan yang ini disampaikanoleh pembuat objek atau karya seni kepada orang lain. Biasanya temasuatu karya akan dipengaruhi oleh banyak faktor, misalnya letak geografis,adat istiadat, budaya, dan lainnya.

	\item Unsur Motif Hias
	\item[] Motif hias adalah pola atau gambar yang menjadi hiasan padasuatu objek atau produk. Tujuan menambahkan motif hias pada suatuobjek adalah untuk menambah nilai keindahan/ estetika pada objek atauproduk tersebut.
\end{enumerate}


\section{Arsitektur}%
\label{sec:Arsitektur}

Arsitektur selalu dikaitkan dengan rumah dan bangunan. Karenanya orangawam mengenal istilah arsitektur selalu terkait dengan pembuatan bangunan ataugedung. Pandangan tersebut memang tidak salah. Karena menurut bahasanyaarsitektur berasal dari Bahasa Latin architectus, dan dari bahasaYunani: architekton (master pembangun), arkhi (ketua) +tekton (pembangun,tukang kayu). Hanya saja arsitektur sebenarnya lebih luas dari itu.

Menurut Kamus Besar Bahasa Indonesia (KBBI), makna arsitekturadalah: 1) seni dan ilmu merancang serta membuat konstruksi bangunan,jembatan, dsb; 2) metode dan gaya rancangan suatu konstruksi bangunan.Menurut Wikipedia, arsitektur adalah seni dan ilmu dalam merancangbangunan. Dalam artian yang lebih luas, arsitektur mencakup merancang dan membangun keseluruhan lingkungan binaan, mulai dari level makro yaituperencanaan kota, perancangan perkotaan, arsitektur lanskap, hingga ke levelmikro yaitu desain bangunan, desain perabot dan desain produk. Arsitektur jugamerujuk kepada hasil-hasil proses perancangan tersebut.

Selain itu, dari definisi di atas dapat disimpulkan bahwa arsitektur jugatermasuk di dalamnya bicara soal konstruksi dan kekuatan bangunan. Berikutbeberapa pendapat para tokoh dunia mengenai pengertian arsitektur:

\begin{enumerate}
	\item Marcus Pollio Vitruvius (1486)
	\item[] Mendefinisikan arsitektur sebagai kesatuan dari kekuatan/kekokohan (firmitas), keindahan (venustas), dankegunaan/fungsi (utilitas).

	\item Banhart CL. Dan Jess Stein
	\item[] Arsitektur adalah seni dalam mendirikan bangunan termasukdidalamnya segi perencanaan, konstruksi, dan penyelesaian dekorasinya;sifat atau bentuk bangunan; proses membangun; bangunan dan kumpulanbangunan.

	\item Claudil (1979)
	\item[] Menurutnya, arsitektur adalah sesuatu yang bersifat personal,menyenangkan dan memerlukan pengalaman. Arsitektur adalah hasilpersepsi dan penghargaan manusia terhadap ruang dan bentuk. Ada tigapengalaman arsitektur: aspek fisikal, emosional dan kebutuhan intelektual.

	\item Van Romondt
	\item[] Arsitektur adalah ruang tempat hidup manusia dengan bahagia.Ruang berarti menunjuk pada semua ruang yang terjadi karena dibuat olehmanusia atau juga ruang yang terjadi karena proses alam seperti gua,naungan pohon dan lain-lain.
\end{enumerate}

\subsection{Estetika dalam Arsitektur}%
\label{sub:Estetika dalam Arsitektur}
Sama seperti sebuah bahasa, bila tidak ada bahasa, maka pengetahuantidak tertularkan. Dalam arsitektur, estetika adalah sebuah bahasa visual, yangtidak sama dengan beberapa bahasa estetika yang tidak visual, seperti bahasa itusendiri. Estetika dalam arsitektur memiliki banyak sangkut paut dengan segala yang visual seperti permukaan, volume, massa, elemen garis, dan sebagainya,termasuk berbagai orderharmoni, seperti komposisi.

Secara umum estetika dalam arsitektur bangunan merupakan filsafatkeindahan bentuk dan ruang maka pemaparan estetika arsitektur ini diarahkanpada pembahasan arsitektur sebagai ilmu dan seni. Sebagai ilmu khususnya ilmurancang bangunan, arsitektur ditujukan untuk merancang suatu sistem pewadahanatau ruang. Ruang yang dimaksutkan adalah ruang yang dikondisikan untuk suatukegiatan sehingga bentuk ruang tersebut harus mengikuti bentuk kegiatanpenghuninya seperti kata Lois Sullivans, form follow function artinya bentukmengikuti fungsi.

Arsitektur tidak hanya sebatas bentuk dan fungsi, tetapi juga seni.Sehingga arsitektur tidak bisa dibatasi oleh titik, garis dan bidang melainkan bisaberkembang menjadi sebuah keindahan seni yang bersumber dari nilai-nilaibudaya, moral, kehidupan, sejarah dll. Berarsitektur berarti berbahasa manusia,maksudnya adalah suatu rancangan arsitektur harus mampu mengkomunikasikanmengenai fungsi dan wadah dari arsitektur sendiri, hal itu dapat diwujudkandalam suatu keindahan bentuk dan seni. Keindahan adalah pancarankebenaran kata. Intinya estetika dalam arsitektur yang tertuang dalam bangunanharuslah kontekstual dan komunikatif. Kontekstual berarti ia memenuhihakekatnya sebagai wadah dan fungsi, komunikatif berarti ia memenuhi unsurseni dan keindahan sehingga mencerminkan fungsi.

Menurut Prijotomo dalam Sejarah Arsitektur (2009) bahwa arsitekturadalah bangunan yang memiliki nilai estetika (estetis), atau dapat dikatakansebagai bangunan yang elok. Estetika berkaitan dengan keindahan dan keindahanitu dapat dirasakan dengan indra pengelihatan. Secara etimologi estetika berasaldari kata Aesthetikos (Yunani) yang artinya pengamatan denganpengindraan. Sehingga estetika itu bisa dinikmati oleh semua panca indra manusiabukan hanya mata saja.

Estetika meskipun berkaitan dengan 'rasa' saat melihat bangunan jugadapat dibangun melalui aplikasi teori arsitektur. Inilah mengapa estetika patutdibahasakan dan dibahas dalam alat yang bernama komunikasi. Estetika dapatdimengerti dan dikembangkan melalui pemahaman berbagai hal menyangkut teoriestetika, menjadi dasar bagi banyak cabang seni.Namun melihat berbagai dimensiyang hanyalah sebuah media untuk mencoba menjelaskan apa yang disebut indah,namun tidak pernah bisa menjelaskan apa yang sebenarnya terjadi dalam benakseseorang mempengaruhi bagaimana seorang manusia mengapresiasi keindahan,estetika berkaitan dengan sensasi keindahan. Dalam teori tentang estetika, dicobadijelaskan berbagai sisi yang 'tersentuh' oleh keindahan sebuah obyek. Jadi, apayang indah bagi saya belum tentu indah bagi Anda.


\section{Perilaku}%
\label{sec:Perilaku}


\subsection{Perkembangan Ilmu Perilaku dan Lingkungan}%
\label{sub:Perkembangan Ilmu Perilaku dan Lingkungan}


Perilaku dalam arti empiris mempunyai karakter kasatmata tetapi penyebab terjadinya tidak dapat diamati contoh makan, memasak, bekerja, belajar, tidur, istirahat dsb.

\begin{enumerate}
	\item Perilaku mengenal berbagai tingkatan yaitu perilaku sederhana dan stereotip, perilaku kompleks seperti perilaku sosial manusia.
	\item Perilaku bervariasi dengan klasifikasi kognitif, afektif dan psikomotorik yang mengarah pada rasio dan emosi.
	\item Perilaku bisa disadari dan juga bisa tidak disadari.
\end{enumerate}

Ilmu hubungan perilaku-lingkungan bertujuan membentuk suatu hubungan yang saling menunjang antara manusia sebagai individu ataupun kelompok dan lingkungan fisiknya guna meningkatkan kualitas kehidupan melalui kebijakan perencanaan dan perancangan (Moore, 1976).

Egon Brunswick (1903-1955) menyatakan bahwa lingkungan fisik mempengaruhi manusia tanpa manusia sendiri menyadarinya, Brunswick inilah yang pertama kali menggunakan istilah Psikologi lingkungan.

Ilmu perilaku-lingkungan mempunyai karakteristik seperti hubungan perilaku-lingkungan adalah satu unit yang dipelajari dalam keadaan saling terkait tidak berdiri sendiri seperti hubungan kebisingan dengan tingkat konsentrasi seseorang atau kondisi fasilitas kantor dengan produktifitas seseorang.

\begin{enumerate}
	\item Hubungan antara lingkungan dan manusia serta perilakunya adalah hubungan timbal balik, saling terkait, dan saling mempengaruhi.
	\item Misalkan dalam ranah pendidikan, apakah sarana dan prasarana belajar yang menyebabkan mutu dan prestasi belajar disekolah.
\end{enumerate}

Perilaku-lingkungan tidak hanya memusatkan pada masalah teoretis atau terapan namun menitik beratkan pada keduanya.

\begin{enumerate}
	\item Interdisipliner yang berarti karena ruang lingkupnya yang bermacam-macam maka dalam penelitiannya harus bekerjasama dengan berbagai disiplin ilmu.
	\item Dari penjelasan mengenai kajian tersebut hendaknya hasil studi perilaku-lingkungan dapat menjadi panduan desain untuk meningkatkan kualitas desain pada tipe dan tatanan tertentu dan untuk kelompok pengguna tertentu.
\end{enumerate}

\subsection{Tinjauan Teori Arsitektur}%
\label{sub:tinjauan_teori_arsitektur}

Arsitektur adalah kristalisasi dari pandangan hidup sehingga arsitektur bukan semata-mata teknik dan estetika bangunan atau bahkan terpecah-pecah menjadi kelompok seperti ranah keteknikan, ranah seni atau ranah sosial.

\begin{enumerate}
	\item Arsitektur adalah ruang fisik untuk aktivitas manusia yang memungkinkan pergerakan manusia dari satu ruang ke ruang lainnya, yang menciptakan tekanan antara ruang dalam bangunan dan ruang luar.
	\item Arsitektur dapat terbentuk karena ada persepsi dan imajinasi manusia sebagai pengguna.
\end{enumerate}

\subsection{Pendekatan Desain dalam Arsitektur Perilaku}%
\label{sub:pendekatan_desain_dalam_arsitektur_perilaku}

Dari ulasan mengenai fokus teori arsitektur tersebut, perubahan atau perkembangan teori arsitektur dan pengambilan keputusan desain harus mempertimbangkan manusia sebagai suatu entitas spiritual, bukan hanya entitas fisik, agar hasil desain dapat mencapai sasaran yang dituju.

Pendekatan desain dalam arsitektur perilaku dibagi menjadi dua yaitu pendekatan Cybernatics (sibernatik) dan Teori positif:

\subsubsection{Cybernatics}%
\label{ssub:cybernatics}

Cybernatics adalah pendekatan desain lingkungan yang menekankan perlunya mempertimbangkan kualitas lingkungan yang dihayati oleh pengguna dan pengaruhnya bagi pengguna lingkungan tersebut. Pendekatan ini secara holistik mengaitkan berbagai fenomena yang mempengaruhi hubungan antara manusia dan lingkungannya termasuk lingkungan fisik dan sosial.

Dimensi pembentuk desain lingkungan sibernatik adalah * Keinginan klien dengan tiga tingkatan kebutuhan pengguna yaitu tingkat kesehatan atau keselamatan dan keamanan, tingkat fungsi dan efisiensi dan tingkat kepuasan psikologis.

\begin{enumerate}
	\item Elemen kerangka hunian yaitu bangunan atau setting berupa tata perilaku (behavioural setting), Barker (1968) dan archetypal, Spivak (1973).
	\item Penghuni, dibedakan berdasarkan siklus kehidupan misal balita, anak-anak, remaja, orang tua, atau penyandang cacat fisik dan cacat mental.
	\item Kebutuhan lain, seperti kebutuhan budaya dan adat.
\end{enumerate}

\subsubsection{Teori positif}%
\label{ssub:teori_positif}

Teori positif mencakup pengertian tentang lingkungan dan perannya bagi kehidupan manusia, Amos (1994). Teori positif merupakan suatu proses kreatif yang mencakup pembentukan struktur konseptual, baik untuk menata maupun untuk menjelaskan hasil suatu pengamatan.

Tujuannya adalah agar struktur ini dapat digunakan untuk menjelaskan apa yang terjadi dan membuat prediksi mengenai apa yang mungkin terjadi. Nilai dari teori positif ini tergantung pada kekuatan penjelasan dan prediksinya. Teori yang berhasil adalah teori yang sederhana tetapi mampu membantu kita dalam memprediksi dengan akurat.

Contoh dari teori positif dalam perancangan adalah kreativitas desain yang dibuat tidak mengabaikan akibat dari desainnya sendiri. Teori positif bertujuan meningkatkan kesadaran mengenai perilaku mana dalam lingkungan yang penting bagi manusia dalam mengambil keputusan desain.

Dalam proses arsitektur yang kreatif, ada empat dimensi studi perilaku-lingkungan yaitu manusia, perilaku, lingkungan, dan waktu.

\subsection{Proses Individu dalam Perilaku Arsitektur}%
\label{sub:Proses Individu dalam Perilaku Arsitektur}

Proses individual sebagai faktor pengaruh dan yang mempengaruhi perilaku manusia meliputi hal-hal sebagai berikut :

\begin{itemize}
	\item Persepsi lingkungan
	\item Kognisi spasial
	\item Perilaku spasial
\end{itemize}



\begin{enumerate}
	\item Persepsi lingkungan adalah proses bagaimana manusia menerima informasi mengenai lingkungan sekitarnya dan bagaimana informasi mengenai ruang fisik tersebut diorganisasikan ke dalam pikiran manusia.
	\item Kognisi spasial adalah keragaman proses berpikir selanjutnya mengorganisasikan, menyimpan dan mengingat kembali informasi mengenai lokasi, jarak dan tatanan dalam lingkungan fisik.
	\item Perilaku spasial adalah perilaku yang menunjukkan hasil yang termanifestasikan dalam tindakan dan respons seseorang, termasuk deskripsi dan preferensi personal respons emosional ataupun evaluasi kecenderungan perilaku yang muncul dalam interaksi manusia dengan lingkungan fisiknya.
\end{enumerate}

\subsection{Lingkungan}%
\label{sub:lingkungan}

Lingkup lingkungan dibedakan menjadi empat macam yaitu :
\begin{itemize}
	\item Lingkungan fisik terdiri atas terestrial atau latar geografis
	\item Lingkungan sosial terdiri atas organisasi sosial kelompok interpersonal
	\item Lingkungan psikologikal terdiri atas imaji yang dimiliki orang dalam benaknya
	\item Lingkungan behavioural mencakup elemen-elemen yang menjadi pencetus respons seseorang.
\end{itemize}

Perbedaan utama dari klasifikasi tersebut adalah perbedaan antara lingkungan obyektif yang nyata di sekitar seorang individu dan lingkungan fenomenologis yang dihayati (perceived) dan yang secara sadar ataupun tidak mempengaruhi pola perilaku dan emosi seseorang.

Gibson (1966), menjelaskan lingkungan yang potensial bagi perilaku dan efektif terdiri atas segala sesuatu yang menjadi perhatian orang atau yang digunakan.

\begin{enumerate}
	\item Lingkungan terestrial atau lingkungan geografis adalah lingkungan yang hanya merujuk pada lingkungan alam seperti tanah dan proses terjadinya. Lingkungan terestrial terdiri atas komponen padat, cair dan gas. Lingkungan terestrial dapat dikatakan menjadi sumber bagi banyak pengalaman manusia sebagai contoh radiant light, ambient, panas, suara bau dan kontak mekanis.
	\item Lingkungan makhluk hidup dalam suatu sistem sosial terdiri atas sekelompok individu yang saling berinteraksi secara langsung ataupun tidak langsung demi kepentingan tertentu. Lingkungan manusia ini terdiri atas sejumlah sistem sosial dengan ketentuan peran dan perilaku anggotanya.
	\item Lingkungan budaya berkembang karena adanya peluang dari sumber alam, minat dan kompetensi manusia. Lingkungan terdiri dari affordances atau kemanfaatan lingkungan untuk berperilaku, tapi tidak semua affordance dapat digunakan, tergantung pada karakteristik budaya, nilai dan kebutuhan individunya.
	\item Lingkungan binaan merupakan bagian dari lingkungan terestrial dan lingkungan budaya. Arsitektur dalam pembahasan studi perilaku-lingkungan dianggap sebagai bentuk adaptasi manusia terhadap kedua lingkungan tersebut serta bagaimana mereka berinteraksi dengan lingkungan tersebut.
	\item Penilaian lingkungan mengacu pada kesan personal seseorang terhadap setting yang bersangkutan, mengacu pada enam jenis kesan personal yaitu deskripsi, evaluasi, penilaian akan keindahan, reaksi emosional, makna dan sikap kepedulian yang dikembangkan si pengamat terhadap setting tersebut.
\end{enumerate}

\subsection{Persepsi}%
\label{sub:persepsi}

Sebagian besar dari arsitektur dibentuk oleh persepsi manusia. Persepsi yang dimaksud adalah proses memperoleh atau menerima informasi dari lingkungan. Proses ini dibagi menjadi beberapa kelompok pendekatan yaitu :

\begin{enumerate}
	\item Pendekatan konvensional yaitu pendekatan yang didasari sensori atau stimuli yang berhubungan dengan pengindraan, yang kadang ditafsirkan sebagai penafsiran pengalaman.
	\item Pendekatan ekologis adalah pendekatan berdasarkan informasi, seorang individu tidaklah menciptakan makna dari apa yang diindrakannya.
	\item Teori gestalt dalam formal-design mengedepankan fenomena design tiga hal yang menjadi pokok teori ini adalah form, isomorphism dan field forces.
	\item Pendekatan ekologis dalam design adalah pendekatan yang bersifat radikal karena sangat bertolak belakang dengan isomorphism dari teori gestalt. Gibson (1966), tidak menganggap indra sebagai saluran pengindraan tetapi suatu sistem perseptual.
\end{enumerate}

\subsection{Proses Sosial dalam Perilaku Arsitektur}%
\label{sub:proses_sosial_dalam_perilaku_arsitektur}

Proses sosial dalam perilaku arsitektur adalah hal yang dipersepsikan manusia tentang lingkungan berkaitan dengan ruang sekitarnya, baik ruang natural maupun ruang buatan. Aspek sosialnya adalah bagaimana manusia berbagi dan membagi ruang dengan sesamanya.

Dalam memenuhi kebutuhan sosialnya manusia berperilaku sosial didalam lingkungannya yang dapat diamati dari : Fenomena perilaku lingkungan, kelompok pemakai dan tempat terjadinya aktifitas.

Perilaku interpersonal manusia tersebut meliputi :

\begin{itemize}
	\item Ruang personal berupa domain kecil sejauh jangkauan manusia yang dimiliki semua orang.
	\item Teritorialitas yaitu kecenderungan untuk menguasai daerah yang lebih luas bagi penggunaan oleh seseorang atau sekelompok pemakai atau bagi fungsi tersebut.
	\item Kesesakan dan kepadatan adalah keadaan apabila ruang fisik yang tersedia sangat terbatas dibandingkan jumlah penggunanya.
	\item Privasi adalah usaha untuk mengoptimalkan pemenuhan kebutuhan sosial manusia
\end{itemize}


\subsection{Cerminan}%
\label{sub:Cerminan}

\begin{itemize}
	\item Menyebrang tidak pada tempatnya, kegagalan fungsi fasilitas, perencanaan atau perilaku yang membudaya melanggar??
	\item Penjara di austria, dampak perilaku dari pengguna?? Arsitektur sebagai lingkungan binaan berhasilkah??
	\item Masih dalam kaitan penjara di Austria, dalam kondisi seperti ini perilaku apakah yang ditimbulkan??
	\item Pemukiman kumuh bantaran kali, masalah sosial apa yang akan muncul?? Bagaimana alternatif desain pemukiman masyarakat bawah seharusnya??
\end{itemize}

\begin{figure}
	\begin{center}
		\includegraphics[width=.82\textwidth]{jpo.png}
		% filename may not contain space but _
	\end{center}
	\caption{Jembatan Penyebrangan Orang}
	\label{fig:jpo}
\end{figure}

\begin{figure}
	\begin{center}
		\includegraphics[width=.82\textwidth]{justice_center_penjara_mewah_di_austria.jpg}
		% filename may not contain space but _
	\end{center}
	\caption{Penjara Mewah di Austria}
	\label{fig:penjara}
\end{figure}

\begin{figure}
	\begin{center}
		\includegraphics[width=.82\textwidth]{permukiman_kumuh}
		% filename may not contain space but _
	\end{center}
	\caption{Permukiman Kumuh}
	\label{fig:}
\end{figure}




\bibliographystyle{apalike}
\bibliography{../filename.bib} % required for citecompletion, not work in mainfile
\end{document}


